\begin{longtable}{TD}
\caption{Task Behavior Graph terms.}
\label{tab:terms-task-behavior-graphs}\\

\toprule
Term & Definition \\
\midrule
\endfirsthead

\toprule
Term & Definition \\
\midrule
\endhead

\bottomrule
\endfoot

\termrow{Task Behavior Graph}{
A typed graph associated with a task kind or task family. It connects canonical
features across influence domains and supports estimation of resource demand,
worker occupancy, uncertainty, and operational risk. A Task Behavior Graph is
not assumed to prove causality by itself; it makes modeled influence
assumptions explicit.
}

\termrow{graph node}{
An element of a Task Behavior Graph representing a canonical feature, output,
dimension, intermediate signal, or model concept. Each node belongs to one or
more influence domains.
}

\termrow{graph edge}{
A modeled influence relationship between two graph nodes. A graph edge is not
a proven causal link unless validated by additional evidence.
}

\termrow{influence role}{
The semantic role of a node or edge in the graph, such as baseline demand,
contextual amplification, additive overhead, symptom, outcome link, risk
propagation, uncertainty indicator, or profiling trigger.
}

\termrow{baseline demand edge}{
An edge that connects an intrinsic feature to a baseline resource-demand or
occupancy estimate. For example, payload size may contribute to network I/O or
disk I/O baseline demand.
}

\termrow{contextual amplification edge}{
An edge by which an execution-context factor amplifies or suppresses demand,
occupancy, uncertainty, or risk.
}

\termrow{additive overhead edge}{
An edge that represents additional fixed or variable overhead, such as setup,
serialization, acknowledgement, commit, or retry bookkeeping.
}

\termrow{symptom edge}{
An edge that connects an influence to a symptom or observed effect, such as
network degradation contributing to retries or longer wall-clock duration.
}

\termrow{outcome link}{
An edge that connects modeled behavior to an outcome, such as timeout, failure,
success, cancellation, or dead-lettering.
}

\termrow{risk propagation edge}{
An edge that connects demand, occupancy, uncertainty, or outcome to an
operational-risk dimension.
}

\termrow{uncertainty indicator}{
A node or edge role indicating that a signal changes uncertainty or confidence
rather than directly changing resource demand.
}

\termrow{sensitivity}{
The estimated strength of a relationship between graph nodes. Sensitivity may
be expert-defined, learned, calibrated, or represented as a qualitative level.
}

\termrow{relevance}{
Whether a feature or relationship matters for a task kind, task family, or
behavior profile. Relevance is distinct from sensitivity.
}

\termrow{graph evaluation}{
The process of applying a Task Behavior Graph to canonical features and
execution context in order to estimate demand, occupancy, uncertainty,
residuals, and risk.
}

\termrow{graph template}{
A reusable graph structure that can be specialized for multiple related task
kinds within a task family.
}

\termrow{model assumption}{
An explicit assumption encoded in the graph, feature mapping, or policy. Model
assumptions should be validated, calibrated, or revised when observations
contradict them.
}

\end{longtable}
